\begin{titlepage}
\thispagestyle{empty}
\vspace*{2cm}
\noindent{\color{apagado}\large\textsc{Psicoestadística Inferencial}}\par
\vspace{0.4em}
\noindent{\Huge\bfseries Unidad 2 · Probabilidad}\par
\vspace{0.3em}
\noindent{\LARGE Cuaderno de ejercicios}\par
\vspace{1.2em}
\noindent\rule{\linewidth}{1pt}\par
\vspace{1.2em}

\noindent Los ejercicios están ordenados por dificultad dentro de cada
apartado, y arrancan deliberadamente por lo más elemental. Si el primero de
un apartado te parece demasiado fácil, está bien: existe para que nadie quede
afuera antes de empezar.\par
\vspace{1em}

\noindent\begin{tabular}{@{}p{2.6cm}p{10.5cm}@{}}
\toprule
\textbf{Nivel} & \textbf{Qué se entrena} \\
\midrule
\textcolor{nivelcero}{\textbf{N0 · Reconocer}}   & Vocabulario y distinciones. Sin cálculo. \\[0.3em]
\textcolor{niveluno}{\textbf{N1 · Contar}}       & Identificar qué va arriba y qué abajo en la fracción. \\[0.3em]
\textcolor{niveldos}{\textbf{N2 · Una fórmula}}  & Aplicarla una vez, con números chicos. \\[0.3em]
\textcolor{niveltres}{\textbf{N3 · Datos reales}}& Sobre las 200 fichas del servicio de salud mental. \\[0.3em]
\textcolor{nivelcuatro}{\textbf{N4 · Decidir}}   & Elegir la herramienta, detectar el error, interpretar. \\
\bottomrule
\end{tabular}\par

\vspace{1.6em}
\noindent Los ejercicios de nivel 3 y 4 usan el mismo archivo de 200 fichas
que la herramienta interactiva de la materia, así que los números coinciden
con los que se ven en pantalla.\par

\vfill
\noindent{\color{apagado}\small
\ifdocente\textbf{Versión del docente} — incluye la solución paso a paso de
cada ejercicio.\else\textbf{Versión del estudiante} — las respuestas están al
final; el desarrollo se trabaja en clase.\fi}\par
\vspace{0.4em}
\noindent{\color{apagado}\small Ronald Martínez Jiménez}\par
\end{titlepage}
